\documentclass[12pt]{amsart}

\usepackage{amsmath,amssymb,amsthm}
\usepackage{mathtools}
\usepackage{booktabs}
\usepackage[margin=1.1in]{geometry}
\usepackage[
  pdftitle={A shifted indefinite theta completion over Q(sqrt(5)): an asymmetric boundary derivative and a row-model factorization},
  pdfauthor={Xiang Huang}
]{hyperref}

\newtheorem{theorem}{Theorem}[section]
\newtheorem{proposition}[theorem]{Proposition}
\newtheorem{lemma}[theorem]{Lemma}
\newtheorem{corollary}[theorem]{Corollary}
\newtheorem{conjecture}[theorem]{Conjecture}
\theoremstyle{definition}
\newtheorem{definition}[theorem]{Definition}
\newtheorem{example}[theorem]{Example}
\theoremstyle{remark}
\newtheorem{remark}[theorem]{Remark}

\DeclareMathOperator{\N}{N}
\DeclareMathOperator{\sgn}{sgn}
\newcommand{\ZZ}{\mathbb{Z}}
\newcommand{\QQ}{\mathbb{Q}}
\newcommand{\RR}{\mathbb{R}}
\newcommand{\OK}{\mathcal{O}_K}
\newcommand{\eps}{\varepsilon}
\newcommand{\ph}{\varphi}
\newcommand{\Acone}{\mathcal{A}}
\newcommand{\Dcone}{\mathcal{D}}
\newcommand{\MK}{\mathrm{MK}}
\newcommand{\Splus}{\mathcal{S}^{+}}
\newcommand{\Sminus}{\mathcal{S}^{-}}

\title[A shifted indefinite theta completion over $\QQ(\sqrt5)$]
{A shifted indefinite theta completion over $\QQ(\sqrt5)$:\\
an asymmetric boundary derivative and a row-model factorization}
\author{Xiang Huang}
\date{}

\begin{document}

\begin{abstract}
We study a norm-supported cone series $B(q)$ over $\QQ(\sqrt5)$ and
two structures attached to it.  First, a coefficientwise row model
motivated by the even-row missing kernel in Chan's $\Theta_{10}$
dissection factors as $-q^{-1}j(q;q^5)^2B(q)$.  Second, the normalized
series $F(\tau)=q^{1/10}B(q)$ is the holomorphic sign kernel of an
explicit rank-two indefinite theta datum in the sense of Zwegers.  We
write the exact characteristic vectors and the resulting componentwise
weight-one transformation formulas.  One boundary correction vanishes by
an odd-theta pairing.  For the other boundary we give an index-four
orthogonal decomposition and compute
$\partial_{\bar\tau}\widehat F$ as a nonzero sum of a holomorphic
weight-$\tfrac12$ theta factor times an antiholomorphic unary
weight-$\tfrac32$ theta factor.  The corresponding $\xi_1$ image is
nonholomorphic, so this natural completion is not harmonic and does
not by itself establish ordinary weight-one mock modularity in the
harmonic-Maass sense.  The result is instead an explicit shifted indefinite theta
completion with a mixed theta derivative.  The arithmetic theory of
$B$ is independent and belongs to a separate companion paper.
\end{abstract}

\maketitle

%----------------------------------------------------------------------
\section{Introduction}\label{sec:intro}
%----------------------------------------------------------------------

The identities of Andrews, Dyson, and Hickerson and of Cohen connect
$q$-series, real quadratic fields, and Maass waveforms
\cite{jADH88,jCo88}.  Their coefficients admit an ideal-theoretic
interpretation and display the multiplicativity expected from a Hecke
character.  Related real-quadratic constructions were developed in
\cite{jCFLZ04,jBrKa11,jLoOs15}.  The series studied here arises from a
different source: Chan's dissection of the tenth-order theta function
$\Theta_{10}$ \cite{jCh10}.

The companion arithmetic paper \cite{companion} studies the independently
defined cone series
\[
  B(q)=\sum_{N\geq 0} B_N q^N
\]
and proves the finite lattice, norm, and coefficient identities needed
to define it.  This paper studies the analytic nature of that already
defined series.  Its main point is that the natural completion has one
vanishing boundary correction and one surviving boundary correction.
The asymmetry is visible both in the unequal norms
$-5$ and $-20$ of the two negative vectors and in the parity of the
transverse theta sums.

Our principal analytic result is the following.  The theta components
$\theta_j$ and $g_j$ are defined explicitly in
Section~\ref{sec:shadow}.

\begin{theorem}[Exact completion and differential image]\label{thm:main}
Let $F(\tau):=q^{1/10}B(q)$, where $q=e^{2\pi i\tau}$.  For
\[
  A=\begin{pmatrix}1&0\\0&-5\end{pmatrix},\qquad
  a=\left(\frac12,\frac1{10}\right),\qquad
  b=\left(\frac12,-\frac1{10}\right),
\]
and the negative-vector pair $c_2=(-5,3)$, $c_1=(0,1)$, let
$\vartheta_{a,b}^{c_2,c_1}$ denote Zwegers' completed indefinite theta
function for $Q_0(v)=\tfrac12v^tAv$.  Then
\[
  \widehat F(\tau)
  :=\frac12e^{-3\pi i/5}\vartheta_{a,b}^{c_2,c_1}(\tau)
\]
has holomorphic sign part $F$ and satisfies
\[
  \widehat F(\tau+1)=e^{\pi i/5}\widehat F(\tau),
\]
\[
  \widehat F(-1/\tau)
  =\frac{\tau}{2\sqrt5}\sum_{m=0}^{4}
    \vartheta_{\,b+(0,m/5),\,-a}^{c_2,c_1}(\tau).
\]
Writing $Y=\operatorname{Im}(\tau)$, its differential images are
\[
  \frac{\partial\widehat F}{\partial\overline\tau}
  =-\frac{i}{4\sqrt{10Y}}
    \sum_{j=0}^{3}(-1)^j\theta_j(\tau)\overline{g_j(\tau)},
\]
and, for $\xi_1(f)=2iY\,\overline{\partial f/\partial\overline\tau}$,
\[
  \xi_1(\widehat F)
  =-\frac{\sqrt Y}{2\sqrt{10}}
    \sum_{j=0}^{3}(-1)^j\overline{\theta_j(\tau)}g_j(\tau).
\]
These expressions are nonzero and the $\xi_1$ image is not
holomorphic.  Hence $\widehat F$ is a component of a finite-dimensional
vector-valued weight-one real-analytic modular completion, but is not a
weight-one harmonic Maass form.
\end{theorem}

Here
\[
  j(z;q):=(z;q)_\infty(q/z;q)_\infty(q;q)_\infty
\]
is Jacobi's product.  We do not attach a stronger mock-modular label
to the residual product here: doing so requires carrying the exact
monomial normalization and the vector-valued representation through
the Jacobi factor.  The formal product identity is stated independently
in Theorem~\ref{thm:mk-factorization}.

\subsection*{Relation to earlier work}
Zwegers' thesis supplies the general indefinite-theta completion and
its characteristic transformation formulas \cite{oZw02}.  His later
paper on mock Maass theta functions places ADH- and Cohen-type sums in
a broader real-analytic setting \cite{jZw12}.  The contribution here is
the exact specialization for this shifted $\QQ(\sqrt5)$ cone: the
characteristics, the asymmetric boundary decomposition, and the mixed
theta derivative.  Hickerson--Mortenson's treatment of Hecke-type
double sums gives complementary $q$-series background
\cite{jHiMo14}.  We do not infer a harmonic-Maass or ordinary mock
classification merely from the existence of a Zwegers completion.

\subsection*{Organization}
Section~\ref{sec:foundations} recalls only the definitions needed from
the companion paper.  Section~\ref{sec:arithmetic-recall} gives a brief
arithmetic separation statement.  Section~\ref{sec:slab} proves the
slab identity, Section~\ref{sec:pell} gives the corrected Pell
coordinates and sign calculation, and Section~\ref{sec:shadow}
computes the completion and mixed boundary derivative.  Section~\ref{sec:comparison}
compares the construction with the Maass cases, and
Section~\ref{sec:status} records computation, formalization scope, and
open questions.

%----------------------------------------------------------------------
\section{The residual series recalled from the companion paper}
\label{sec:foundations}
%----------------------------------------------------------------------

We keep the common foundation short and refer to \cite{companion} for
the lattice bijection, norm interpretation, unit action, prime
classification, and their proofs.

\begin{definition}\label{def:E}
For $k,r\in\ZZ$, put
\[
  H(k,r):=4k^2+2k+r^2+(6k+1)r,
  \qquad E(k,r):=\frac{H(k,r)}2.
\]
The integer $H(k,r)$ is even.  Define the two quadrant cones
\[
  \Acone:=\{(k,r)\in\ZZ^2:k\geq0,\ r\geq0\},
  \qquad
  \Dcone:=\{(k,r)\in\ZZ^2:k<0,\ r<0\}.
\]
On these two cones $E(k,r)\geq0$.
\end{definition}

\begin{definition}\label{def:B}
The coefficient and generating series are
\[
  B_N:=
  \sum_{\substack{(k,r)\in\Dcone\\E(k,r)=N}}(-1)^r
  -
  \sum_{\substack{(k,r)\in\Acone\\E(k,r)=N}}(-1)^r,
  \qquad
  B(q):=\sum_{N\geq0}B_Nq^N.
\]
Each coefficient is a finite signed sum.
\end{definition}

The companion paper proves that the map
\[
  \beta(k,r)=(r-2k)+(4k+3r+1)\ph
\]
is a bijection from $\ZZ^2$ to a translate of an index-$10$ sublattice
of $\ZZ[\ph]$ and that
\[
  -\N(\beta(k,r))=10E(k,r)+1.
\]
Section~\ref{sec:mk-factorization} defines a coefficientwise row model
motivated by an even-row missing-kernel expression in the $\Theta_{10}$
project and proves that the model factors exactly as
\[
  \mathcal R(q)=-q^{-1}j(q;q^5)^2B(q)
\]
in the $E$-level exponent chart
(Theorem~\ref{thm:mk-factorization}; an earlier project record
displayed this identity with a mixed normalization --- see
Remark~\ref{rem:normalization}).  The current formal theorem is about
the row model defined here; an identification with the older source
object \texttt{rhoEvenSourceOldRow\_missing} would require a separate
source-to-model theorem and is not claimed.  The finite algebraic
results above are inputs here, not reproved in this paper.

\begin{remark}[Hickerson--Mortenson form]
In the notation of \cite{jHiMo14}, $B$ coincides with the
Hecke-type double sum
\[
  -f_{1,3,4}(X,\,-X^3,\,X);
\]
this follows exactly by termwise comparison of the two cones, signs,
and exponents.  The coefficient implementation was additionally
checked through $N=200$.
\end{remark}

%----------------------------------------------------------------------
\section{Arithmetic separation in five sentences}
\label{sec:arithmetic-recall}
%----------------------------------------------------------------------

The companion paper identifies the defining translate in
$\ZZ[\ph]$ and proves the norm formula above \cite{companion}.
The fundamental totally positive unit does not stabilize that translate
at one step, so the element sum does not descend by the classical ADH
route to a single ideal character.
The resulting coefficient sequence is nonmultiplicative, with
$B_{34}=3$ but $B_1B_3=-2$.
This witness rules out the classical one-dimensional Hecke-character
explanation; it does not rule out every automorphic interpretation.
The prime-value classification and all sector-residue questions belong
to the companion paper and are not duplicated here \cite{companion}.

%----------------------------------------------------------------------
\section{The involution and the exact slab identity}\label{sec:slab}
%----------------------------------------------------------------------

The analytic form of $B$ becomes visible after a sign-reversing
involution removes the paired interior.

\begin{definition}\label{def:sigma}
For fixed $k\in\ZZ$, define
\[
  \sigma_k(r):=-r-6k-1.
\]
Define the two slabs
\[
  \Splus:=\{(k,r)\in\ZZ^2:k\geq0,\ r\leq-6k-1\},
\]
\[
  \Sminus:=\{(k,r)\in\ZZ^2:k\leq-1,\ r\geq-6k\}.
\]
\end{definition}

\begin{lemma}[The involution]\label{lem:sigma}
For every $k,r\in\ZZ$,
\[
  \sigma_k(\sigma_k(r))=r,
  \qquad
  E(k,\sigma_k(r))=E(k,r),
  \qquad
  (-1)^{\sigma_k(r)}=-(-1)^r.
\]
The map $\sigma_k$ has no fixed point in $\ZZ$.
\end{lemma}

\begin{proof}
Write $C=6k+1$.  Since
\[
  H(k,r)=2k(2k+1)+r(r+C),
\]
we have
\[
  (-C-r)(-r)=r(r+C),
\]
which proves preservation of $H$ and hence of $E$.  The number $C$ is
odd, so $(-1)^{-C-r}=-(-1)^r$.  A fixed point would give
$2r=-C$, impossible because the right side is odd.
\end{proof}

\begin{theorem}[Slab identity]\label{thm:slab}
Coefficientwise as a formal power series,
\[
  B(q)=
  \sum_{(k,r)\in\Splus}(-1)^r q^{E(k,r)}
  -
  \sum_{(k,r)\in\Sminus}(-1)^r q^{E(k,r)}.
\]
\end{theorem}

\begin{proof}
For $(k,r)\in\Acone$, the point
$(k,\sigma_k(r))$ lies in $\Splus$, and this assignment is a
bijection from $\Acone$ to $\Splus$.  Lemma~\ref{lem:sigma} gives
\[
  -(-1)^r q^{E(k,r)}
  =(-1)^{\sigma_k(r)}q^{E(k,\sigma_k(r))}.
\]
For $(k,r)\in\Dcone$, the point
$(k,\sigma_k(r))$ lies in $\Sminus$, and the same map is a bijection
from $\Dcone$ to $\Sminus$.  In this case
\[
  (+1)(-1)^r q^{E(k,r)}
  =-(-1)^{\sigma_k(r)}q^{E(k,\sigma_k(r))}.
\]
Adding the two identities proves the result coefficientwise.  No
analytic rearrangement is used.
\end{proof}

\section{A row-model kernel and its exact factorization}
\label{sec:mk-factorization}

Throughout, all series are formal Laurent series over $\ZZ$: a series
is \emph{admissible} if its exponents are bounded below and each
coefficient is a finite integer sum.  Recall from
Definition~\ref{def:E} the quadratics
\[
  H(k,r)=4k^2+2k+r^2+(6k+1)r,\qquad E(k,r)=\tfrac12 H(k,r),
\]
the cones $\Acone=\{k\ge0,\ r\ge0\}$, $\Dcone=\{k<0,\ r<0\}$, the
coefficients $B_N$ of Definition~\ref{def:B}, and the involution
$\sigma_k(r)=-r-6k-1$ of Lemma~\ref{lem:sigma}.  Set further
\[
  Q_{\mathrm{out}}(u,v):=5u^2-3u+5v^2-7v .
\]
Both $Q_{\mathrm{out}}$ and $H$ take only even values, and
$Q_{\mathrm{out}}(u,v)\ge-2$ for all $u,v\in\ZZ$.

\subsection{The row-model kernel}

Motivated by the even-row expression in the $\Theta_{10}$
formalization, consider the four-fold row model with weight
$(-1)^{u+v+r}$, exponent
$Q_{\mathrm{out}}(u,v)+H(k,r)$ (in the dissection chart the exponent
carries an overall factor $9$, which we normalize away), and $(k,r)$
restricted to the two mixed quadrants, with a relative sign:
\[
  R_k:=\begin{cases}\{r\le-1\},&k\ge0,\\[2pt] \{r\ge0\},&k\le-1,\end{cases}
  \qquad
  \varepsilon_k:=\begin{cases}+1,&k\ge0,\\[2pt] -1,&k\le-1.\end{cases}
\]

The four-fold sum must be grouped correctly to be meaningful.  For
fixed $H$-value $m$ the mixed quadrants may contain \emph{infinitely
many} lattice points: the value $m=-6$, for instance, is attained on
$\{k\ge0,\ r\le-1\}$ at
\[
  (1,-3),\ (1,-4),\ (3,-3),\ (3,-16),\ (31,-24),\ (31,-163),\ \dots,
\]
an infinite Pell-automorph orbit family; moreover
$H(k,-3k)=-5k^2-k\to-\infty$, so $H$ is unbounded below there.  Thus
the naive sum over $\ZZ^2_{u,v}\times(\text{mixed quadrants})$ is not
absolutely summable coefficientwise.  The computational probe
\texttt{even\_missing\_kernel\_term} and the Lean declaration
\texttt{evenMissingKernelTerm} use the same row convention; the older
Lean source object is \texttt{rhoEvenSourceOldRow\_missing}.  These
definitions sum \emph{row by row}: for each fixed $k$ the inner
variable $r$ is resolved exactly.  We adopt the same convention.

\begin{definition}[Row-model kernel, row convention]\label{def:mk}
For $k\in\ZZ$ define the row series
\[
  \rho_k(x):=\varepsilon_k
  \sum_{u\in\ZZ}\sum_{v\in\ZZ}\sum_{r\in R_k}
  (-1)^{u+v+r}\,x^{\,Q_{\mathrm{out}}(u,v)+H(k,r)} ,
\]
and the row-model kernel
\[
  \MK(x):=\sum_{k\in\ZZ}\rho_k(x).
\]
Lemma~\ref{lem:mk-welldef} below shows each $\rho_k$ is admissible and
the $k$-sum is locally finite, so $\MK$ is a well-defined Laurent
series with integer coefficients.  All its exponents are even.
\end{definition}

\subsection{Row transport}

The engine of the proof is a row-local transport identity: on each
row, the mixed-quadrant sum equals \emph{minus} the corresponding cone
sum.  It refines the slab identity (Theorem~\ref{thm:slab}) and is
proved by the same involution.

\begin{lemma}[Row transport]\label{lem:row-transport}
Fix $k\in\ZZ$ and $m\in\ZZ$.  All four sums below have finitely many
terms, and
\[
  \sum_{\substack{r\le-1\\ H(k,r)=m}}(-1)^r
  \;=\;-\!\!\sum_{\substack{r\ge0\\ H(k,r)=m}}(-1)^r
  \qquad(k\ge0),
\]
\[
  \sum_{\substack{r\ge0\\ H(k,r)=m}}(-1)^r
  \;=\;-\!\!\sum_{\substack{r\le-1\\ H(k,r)=m}}(-1)^r
  \qquad(k\le-1).
\]
Equivalently, as admissible Laurent series in $x$,
\begin{equation}\label{eq:row-transport}
  \sum_{r\in R_k}(-1)^r x^{H(k,r)}
  \;=\;
  -\sum_{r\in R_k^{c}}(-1)^r x^{H(k,r)},
\end{equation}
where $R_k^{c}=\ZZ\setminus R_k$ is the cone side of the row
($\,r\ge0$ if $k\ge0$; $r\le-1$ if $k\le-1$).
\end{lemma}

\begin{proof}
Finiteness: for fixed $k$, $H(k,\cdot)$ is a monic quadratic in $r$,
so each value is attained at most twice and $H(k,r)\to\infty$ as
$|r|\to\infty$; in particular each row series is admissible.

Let $k\ge0$ and recall $\sigma_k(r)=-r-6k-1$, which preserves
$H(k,\cdot)$ and flips $(-1)^r$ (Lemma~\ref{lem:sigma}).  Split
\[
  \{r\le-1\}=S\sqcup T,\qquad
  S:=\{-6k\le r\le-1\},\qquad
  T:=\{r\le-6k-1\}.
\]
The strip $S$ (empty when $k=0$) is $\sigma_k$-stable: $r\in S$ gives
$\sigma_k(r)=-r-6k-1\in S$.  Since $\sigma_k$ has no fixed points, it
partitions the finite set $S$ into pairs with equal $H$-value and
opposite sign, so $S$ contributes $0$ to every coefficient.  Next,
$r\ge0\iff\sigma_k(r)\le-6k-1$, so $\sigma_k$ is a bijection
$\{r\ge0\}\to T$ preserving $H$ and flipping the sign; hence the
$T$-sum equals minus the $\{r\ge0\}$-sum, coefficient by coefficient.

For $k\le-1$ the argument is the mirror image with
$S=\{0\le r\le-6k-1\}$ (a $\sigma_k$-stable strip of even cardinality
$-6k$) and $T=\{r\ge-6k\}=\sigma_k(\{r\le-1\})$.
\end{proof}

\begin{remark}
Summing \eqref{eq:row-transport} over $k\ge0$ and $k\le-1$ separately
recovers the $\sigma$-straddling description of the row model and,
in the $E$-variable, the slab identity of Theorem~\ref{thm:slab}.  The
row-local form is what is needed here, because only after transport do
the row supports become bounded below uniformly enough to sum over
$k$.
\end{remark}

\subsection{Well-definedness and separation}

\begin{lemma}[Well-definedness]\label{lem:mk-welldef}
Each $\rho_k$ is admissible, and
\begin{equation}\label{eq:row-factored}
  \rho_k(x)=
  \begin{cases}
    -\,\Theta_u(x)\,\Theta_v(x)\displaystyle\sum_{r\ge0}(-1)^r x^{H(k,r)},
    & k\ge0,\\[10pt]
    \;\;\,\Theta_u(x)\,\Theta_v(x)\displaystyle\sum_{r\le-1}(-1)^r x^{H(k,r)},
    & k\le-1,
  \end{cases}
\end{equation}
where
\[
  \Theta_u(x):=\sum_{u\in\ZZ}(-1)^u x^{5u^2-3u},
  \qquad
  \Theta_v(x):=\sum_{v\in\ZZ}(-1)^v x^{5v^2-7v}.
\]
The support of $\rho_k$ is contained in
$[\,4k^2+2k-2,\ \infty)$ for $k\ge0$ and in $[\,4k^2-4k-2,\ \infty)$
for $k\le-1$.  Consequently, for each exponent $n$ only the rows with
$4k^2-4|k|-2\le n$ contribute, and $\MK=\sum_k\rho_k$ is a
well-defined Laurent series.
\end{lemma}

\begin{proof}
Fix $k$ and an exponent $n$.  A term of $\rho_k$ contributes to $x^n$
only if $Q_{\mathrm{out}}(u,v)=a$ and $H(k,r)=m$ with $a+m=n$.  Since
$a\ge-2$ and $m\ge\min_{r\in R_k}H(k,r)>-\infty$, only finitely many
splittings $(a,m)$ occur; each $a$ is attained by finitely many
$(u,v)$ (the quadratic $Q_{\mathrm{out}}$ is positive definite up to
the bounded linear part: $5u^2-3u\le a+2$ and $5v^2-7v\le a+2$ bound
$u,v$), and each $m$ by at most two $r\in R_k$.  Hence $\rho_k$ is
admissible, and by the distributive law for admissible series,
\[
  \rho_k(x)
  =\varepsilon_k\,
   \Bigl(\sum_u(-1)^u x^{5u^2-3u}\Bigr)
   \Bigl(\sum_v(-1)^v x^{5v^2-7v}\Bigr)
   \Bigl(\sum_{r\in R_k}(-1)^r x^{H(k,r)}\Bigr).
\]
Here the index set is the full product $\ZZ_u\times\ZZ_v\times R_k$:
the region constraint involves $(k,r)$ only, the exponent is additively
separated, and the weight $(-1)^{u+v+r}$ is multiplicatively
separated — this is the entire content of the separation step.
Applying Lemma~\ref{lem:row-transport} to the last factor gives
\eqref{eq:row-factored}.

Supports: for $k\ge0$ the map $r\mapsto H(k,r)$ is increasing on
$r\ge0$, so $\min_{r\ge0}H(k,r)=H(k,0)=4k^2+2k$; for $k\le-1$ the
vertex $-(6k+1)/2$ of $H(k,\cdot)$ is positive, so on $r\le-1$ the
minimum is $H(k,-1)=4k^2-4k$.  Since $\Theta_u\Theta_v$ has exponents
$\ge-2$, the stated supports follow from \eqref{eq:row-factored}, and
local finiteness of $\sum_k\rho_k$ is immediate.
\end{proof}

\begin{proposition}[Separation]\label{prop:separation}
As Laurent series,
\[
  \MK(x)=\Theta_u(x)\,\Theta_v(x)\,\bigl(D(x)-A(x)\bigr)
        =\Theta_u(x)\,\Theta_v(x)\,B(x^2),
\]
where $A(x)=\sum_{(k,r)\in\Acone}(-1)^r x^{H(k,r)}$ and
$D(x)=\sum_{(k,r)\in\Dcone}(-1)^r x^{H(k,r)}$.
\end{proposition}

\begin{proof}
Summing \eqref{eq:row-factored} over $k$: the sums
$\sum_{k\ge0}\sum_{r\ge0}$ and $\sum_{k\le-1}\sum_{r\le-1}$ are the
cone series $A(x)$ and $D(x)$, which are admissible with locally
finite row decompositions (row supports $[4k^2+2k,\infty)$ resp.\
$[4k^2-4k,\infty)$ as in Lemma~\ref{lem:mk-welldef}), so
multiplication by the admissible series $\Theta_u\Theta_v$ commutes
with the $k$-summation coefficientwise:
\[
  \MK=\Theta_u\Theta_v\Bigl(\sum_{k\le-1}\ \sum_{r\le-1}(-1)^r x^{H}
      -\sum_{k\ge0}\ \sum_{r\ge0}(-1)^r x^{H}\Bigr)
     =\Theta_u\Theta_v\,(D-A).
\]
Finally $H(k,r)=2E(k,r)$, so $A(x)=\sum_N A_N x^{2N}$ and
$D(x)=\sum_N D_N x^{2N}$ with $A_N,D_N$ the cone counts of
Definition~\ref{def:B}; hence $D(x)-A(x)=B(x^2)$.
\end{proof}

\subsection{The theta legs}

\begin{lemma}[Jacobi legs]\label{lem:jtp}
With $j(z;q):=\sum_{n\in\ZZ}(-1)^n q^{\binom n2}z^n
=(z;q)_\infty(q/z;q)_\infty(q;q)_\infty$,
\begin{enumerate}
\item $\Theta_v(x)=-x^{-2}\,\Theta_u(x)$;
\item $\Theta_u(x)=j(x^2;x^{10})
      =(x^2;x^{10})_\infty(x^8;x^{10})_\infty(x^{10};x^{10})_\infty$;
\item $\Theta_u(x)\,\Theta_v(x)=-x^{-2}\,j(x^2;x^{10})^2$.
\end{enumerate}
\end{lemma}

\begin{proof}
(i) Substitute $u=1-v$ in $\Theta_u$: since
$5(1-v)^2-3(1-v)=5v^2-7v+2$ and $(-1)^{1-v}=-(-1)^v$, the bilateral
sum gives $\Theta_u(x)=-x^{2}\Theta_v(x)$.
(ii) $5u^2-3u=10\binom u2+2u$, so
$\Theta_u(x)=\sum_u(-1)^u (x^{10})^{\binom u2}(x^{2})^{u}
=j(x^2;x^{10})$, and the Jacobi triple product evaluates it as
stated.
(iii) Multiply (i) and (ii).
\end{proof}

\subsection{Main theorem}

\begin{theorem}[Row-model factorization]\label{thm:mk-factorization}
As formal Laurent series,
\[
  \MK(x)\;=\;-x^{-2}\,j(x^2;x^{10})^{2}\,B(x^{2}).
\]
Equivalently, since every exponent of $\MK$ is even, writing
$\mathcal R(q):=\sum_{N}\bigl[x^{2N}\bigr]\MK(x)\,q^{N}$ for the
row model in the halved (that is, $E$-level) variable,
\[
  \boxed{\ \mathcal R(q)\;=\;-q^{-1}\,j(q;q^5)^{2}\,B(q)\ }
\]
with $j(q;q^5)=(q;q^5)_\infty(q^4;q^5)_\infty(q^5;q^5)_\infty$.
\end{theorem}

\begin{proof}
Combine Proposition~\ref{prop:separation} with
Lemma~\ref{lem:jtp}(iii); then substitute $q=x^2$, using
$j(x^2;x^{10})=j(q;q^5)\big|_{q=x^2}$.
\end{proof}

\begin{remark}[Normalization of the previously displayed identity]
\label{rem:normalization}
Earlier project records displayed the candidate identity as
$\mathcal R(q)=-q^{-1}j(q;q^5)^2\,B(q^2)$.  That display mixes the two
exponent charts: the theta legs are written at $E$-level (base $q^5$)
while the cone factor is left at $H$-level.  The internally consistent
forms are the two displayed in
Theorem~\ref{thm:mk-factorization}; the hybrid form fails already at
the constant term ($-3$ versus $-2$).  The coefficientwise
verification (now to $N=2048$) confirms the corrected forms.
\end{remark}

\begin{remark}[Chart with the factor $9$]
In the dissection chart of the formal development the exponent is
$9\bigl(Q_{\mathrm{out}}+H\bigr)=18\,E_{\mathrm{tot}}$.  Thus the
coefficient of $q^{18N}$ there equals the coefficient of $q^N$ in
$-q^{-1}j(q;q^5)^2B(q)$, and every exponent not divisible by $18$ has
coefficient $0$ — in particular the kernel is supported on multiples
of $9$, consistent with the $9\mid e$ phenomenon observed for the
per-fiber counterexamples.
\end{remark}


%----------------------------------------------------------------------
\section{Pell coordinates and the Zwegers cone}\label{sec:pell}
%----------------------------------------------------------------------

Set
\[
  x:=r+3k,\qquad y:=k,
  \qquad P(X,Y):=X^2-5Y^2.
\]
Thus $r=x-3y$ and $k=y$.

\begin{proposition}[Corrected shifted Pell identity]\label{prop:pell-shift}
With
\[
  \mu:=\left(\frac12,\frac1{10}\right),
  \qquad v_{k,r}:=(x,y)+\mu,
\]
we have
\[
  H(k,r)=x^2+x-5y^2-y=P(v_{k,r})-\frac15
\]
and hence
\[
  E(k,r)=\frac12P(v_{k,r})-\frac1{10}.
\]
Moreover $(-1)^r=(-1)^{x+y}$.
\end{proposition}

\begin{proof}
Substituting $r=x-3y$ and $k=y$ gives
\[
  4k^2+6kr+r^2=x^2-5y^2,
  \qquad 2k+r=x-y.
\]
Therefore $H=x^2+x-5y^2-y$.  On the other hand,
\[
  P\left(x+\frac12,y+\frac1{10}\right)
  =x^2+x-5y^2-y+\frac15.
\]
The parity statement follows from $r=x-3y\equiv x+y\pmod2$.
\end{proof}

Let the bilinear form associated with $P$ be
\[
  \mathcal B((X,Y),(X',Y')):=2(XX'-5YY').
\]
Choose the two negative vectors
\[
  c_1=(0,1),\qquad c_2=(-5,3),
\]
for which
\[
  P(c_1)=-5,\qquad P(c_2)=-20.
\]

\begin{proposition}[Corrected bilinear signs]\label{prop:bilinear}
For $v_{k,r}=(x+\tfrac12,y+\tfrac1{10})$,
\[
  \mathcal B(v_{k,r},c_1)=-(10k+1)
\]
and
\[
  \mathcal B(v_{k,r},c_2)=-10(6k+r)-8.
\]
Neither expression vanishes for integral $k,r$.
\end{proposition}

\begin{proof}
Since $y=k$,
\[
  \mathcal B(v_{k,r},c_1)
  =2\left(0-5\left(k+\frac1{10}\right)\right)
  =-(10k+1).
\]
Since $x=r+3k$ and $y=k$,
\begin{align*}
  \mathcal B(v_{k,r},c_2)
  &=2\left[-5\left(x+\frac12\right)
       -15\left(y+\frac1{10}\right)\right]\\
  &=-10x-30y-8\\
  &=-10(r+6k)-8.
\end{align*}
The first expression is odd, and the second is congruent to $2$ modulo
$10$, so neither is zero.
\end{proof}

\begin{theorem}[The slab as a sign-difference cone]
\label{thm:sign-cone}
For every $(k,r)\in\ZZ^2$,
\[
  \mathbf 1_{\Splus}(k,r)-\mathbf 1_{\Sminus}(k,r)
  =\frac12\left(
    \sgn\mathcal B(v_{k,r},c_2)
    -\sgn\mathcal B(v_{k,r},c_1)
  \right).
\]
\end{theorem}

\begin{proof}
By Proposition~\ref{prop:bilinear},
\[
  \sgn\mathcal B(v_{k,r},c_1)=
  \begin{cases}
    -1,&k\geq0,\\
    +1,&k\leq-1,
  \end{cases}
\]
and
\[
  \sgn\mathcal B(v_{k,r},c_2)=
  \begin{cases}
    +1,&6k+r\leq-1,\\
    -1,&6k+r\geq0.
  \end{cases}
\]
The four possible sign pairs give respectively the values $1$, $-1$,
and $0$ off the two slabs, exactly as stated.
\end{proof}

Put $q=e^{2\pi i\tau}$ with $\tau$ in the upper half-plane.  Combining
Theorem~\ref{thm:slab}, Proposition~\ref{prop:pell-shift}, and
Theorem~\ref{thm:sign-cone} gives the exact expression
\begin{align}
  B(q)
  ={}&q^{-1/10}
  \sum_{(x,y)\in\ZZ^2}(-1)^{x+y}
  \frac{\sgn\mathcal B((x,y)+\mu,c_2)
       -\sgn\mathcal B((x,y)+\mu,c_1)}2
  q^{P((x,y)+\mu)/2}.
  \label{eq:zwegers-series}
\end{align}

For the exact comparison with Zwegers' convention, put
\[
  Q_0(X,Y):=\frac12(X^2-5Y^2),
  \qquad
  \mathcal B_0((X,Y),(X',Y')):=XX'-5YY'.
\]
Thus $Q_0(v)=\tfrac12v^tAv$ for
$A=\operatorname{diag}(1,-5)$, and
$\mathcal B=2\mathcal B_0$.  Set
\[
  a:=\mu=\left(\frac12,\frac1{10}\right),
  \qquad
  b:=\left(\frac12,-\frac1{10}\right).
\]
For $v=(x,y)+a$ one has
\begin{equation}\label{eq:characteristic-phase}
  e^{2\pi i\mathcal B_0(v,b)}
  =e^{3\pi i/5}(-1)^{x+y}.
\end{equation}

\begin{theorem}[Exact Zwegers component]\label{thm:zwegers-identification}
Let $F(\tau):=q^{1/10}B(q)$.  In the notation of Definition~2.1 of
\cite{oZw02},
\begin{equation}\label{eq:exact-completion}
  \widehat F(\tau)
  =\frac12e^{-3\pi i/5}
    \vartheta_{a,b}^{c_2,c_1}(\tau).
\end{equation}
Its holomorphic sign part is $F$.  Moreover
\begin{equation}\label{eq:T-transform}
  \widehat F(\tau+1)=e^{\pi i/5}\widehat F(\tau)
\end{equation}
and
\begin{equation}\label{eq:S-transform}
  \widehat F(-1/\tau)
  =\frac{\tau}{2\sqrt5}\sum_{m=0}^{4}
    \vartheta_{\,b+(0,m/5),\,-a}^{c_2,c_1}(\tau).
\end{equation}
Thus $\widehat F$ belongs to the finite characteristic orbit supplied
by Zwegers' weight-one transformation formulas.
\end{theorem}

\begin{proof}
Equation~\eqref{eq:zwegers-series}, the identity
\eqref{eq:characteristic-phase}, and the fact that
$\mathcal B$ and $\mathcal B_0$ have the same signs give
\eqref{eq:exact-completion}.  Both vectors are negative for $Q_0$
and lie in the same negative component,
since
\[
  \mathcal B_0(c_1,c_2)=-15<0;
\]
Proposition~\ref{prop:bilinear} shows that no shifted lattice point lies
on either boundary.  Zwegers' completion theorem applies.

For the $T$-transformation, Corollary~2.9 of \cite{oZw02} gives the
factor $e^{-2\pi iQ_0(a)}=e^{-\pi i/5}$ and replaces $b$ by
$a+b+\tfrac12(1,1)$.  The corollary's second exponential factor is
$1$, because
$\mathcal B_0(A^{-1}A^{\ast},a)=\mathcal B_0((1,1),a)=0$.
The new characteristic differs from $b$ by $(1,3/5)$, an element of
$A^{-1}\ZZ^2$, and the characteristic shift contributes
$e^{2\pi i\mathcal B_0(a,(1,3/5))}=e^{2\pi i/5}$.  Their product is
$e^{\pi i/5}$, proving \eqref{eq:T-transform}.

For $S$, the same corollary sums over
$A^{-1}\ZZ^2/\ZZ^2$, represented by $(0,m/5)$ for $0\leq m<5$.
Here $\det A=-5$ and
$\mathcal B_0(a,b)=3/10$, so its prefactor is
$(\tau/\sqrt5)e^{3\pi i/5}$.  Multiplication by the constant in
\eqref{eq:exact-completion} gives \eqref{eq:S-transform}.
\end{proof}

%----------------------------------------------------------------------
\section{One-sided completion and the mixed boundary derivative}
\label{sec:shadow}
%----------------------------------------------------------------------

We now distinguish the two boundary corrections.  The argument is
concrete and does not infer an analytic category merely from the inequality
$P(c_1)\neq P(c_2)$.

\begin{lemma}[Odd-theta cancellation]\label{lem:odd-theta}
For every $a>0$ and every $\tau$ in the upper half-plane,
\[
  \sum_{n\in\ZZ}(-1)^n
  e^{2\pi i\tau a(n+1/2)^2}=0.
\]
\end{lemma}

\begin{proof}
The involution $n\mapsto-1-n$ preserves $(n+\tfrac12)^2$ and changes
$(-1)^n$ to its negative.  It has no integral fixed point, so the terms
cancel in pairs.
\end{proof}

\begin{proposition}[The $c_1$ correction vanishes]
\label{prop:c1-vanishes}
In the Zwegers completion of Theorem~\ref{thm:zwegers-identification},
the boundary correction attached to $c_1=(0,1)$ is identically zero.
\end{proposition}

\begin{proof}
The vector $e_1=(1,0)$ satisfies
\[
  \mathcal B(c_1,e_1)=2\bigl(0\cdot 1-5\cdot 1\cdot 0\bigr)=0,
  \qquad |\det(c_1,e_1)|=1,
\]
so the lattice splits unimodularly in the $c_1$ and $e_1$ directions.
On each longitudinal fiber the transverse variable is $x$ with shift
$\tfrac12$, and the character restricts as
$(-1)^{x+y}=(-1)^{y}\cdot(-1)^{x}$, i.e.\ to $(-1)^n$ on the
transverse fiber.  The correction therefore
factors into its longitudinal error-function term times the transverse
series in Lemma~\ref{lem:odd-theta}.  The latter vanishes identically,
so the complete $c_1$ correction vanishes.
\end{proof}

For the second boundary, take $e_2=(3,-1)$.  Then
$\mathcal B_0(c_2,e_2)=0$ and $|\det(c_2,e_2)|=4$, so the boundary
lattice has four residue components.  For $v=(x,y)+a$, write
\[
  n:=5x+15y+4,
  \qquad
  T:=3x+5y+2.
\]
The inverse formulas are
\[
  x=\frac{3T-n-2}{4},
  \qquad
  y=\frac{3n-5T-2}{20},
\]
so the integer lattice is the disjoint union, for $j\in\ZZ/4\ZZ$, of
the pairs satisfying
\[
  n\equiv5j+4\pmod{20},
  \qquad
  T\equiv3j+2\pmod4.
\]
Moreover
\begin{equation}\label{eq:orthogonal-Q}
  v=\frac{n}{20}c_2+\frac{T}{4}e_2,
  \qquad
  Q_0(v)=\frac{T^2}{8}-\frac{n^2}{40},
  \qquad
  (-1)^{x+y}=(-1)^j.
\end{equation}

Define the transverse and longitudinal theta series
\begin{equation}\label{eq:theta-g-defs}
  \theta_j(\tau):=
  \sum_{\substack{T\in\ZZ\\T\equiv3j+2\ (4)}}q^{T^2/8},
  \qquad
  g_j(\tau):=
  \sum_{\substack{n\in\ZZ\\n\equiv5j+4\ (20)}}nq^{n^2/40}.
\end{equation}
These are vector-valued unary theta components of weights
$\tfrac12$ and $\tfrac32$, respectively.

\begin{proposition}[Exact boundary derivative]\label{prop:shadow}
Put $Y=\operatorname{Im}(\tau)$.
For the completion \eqref{eq:exact-completion},
\begin{equation}\label{eq:dbar-exact}
  \frac{\partial\widehat F}{\partial\overline\tau}
  =-\frac{i}{4\sqrt{10Y}}
    \sum_{j=0}^{3}(-1)^j\theta_j(\tau)\overline{g_j(\tau)}.
\end{equation}
Consequently
\begin{equation}\label{eq:xi-exact}
  \xi_1(\widehat F)
  =-\frac{\sqrt Y}{2\sqrt{10}}
    \sum_{j=0}^{3}(-1)^j\overline{\theta_j(\tau)}g_j(\tau).
\end{equation}
Both expressions are nonzero, and \eqref{eq:xi-exact} is not
holomorphic.
\end{proposition}

\begin{proof}
For the second boundary,
\[
  \mathcal B_0(c_2,v)=-n,
  \qquad Q_0(c_2)=-10.
\]
Zwegers' error kernel is therefore
$E(-n\sqrt{Y/10})$.  Since
$E'(z)=2e^{-\pi z^2}$,
\[
  \frac{\partial}{\partial\overline\tau}
  E(-n\sqrt{Y/10})
  =-\frac{in}{2\sqrt{10Y}}e^{-\pi n^2Y/10}.
\]
Insert this identity in \eqref{eq:exact-completion}, use
\eqref{eq:orthogonal-Q}, and group by $j$.  The constant phase in
\eqref{eq:characteristic-phase} cancels the prefactor in
\eqref{eq:exact-completion}, yielding \eqref{eq:dbar-exact}.  The
$c_1$ derivative contributes zero by
Proposition~\ref{prop:c1-vanishes}.  Taking complex conjugates and
multiplying by $2iY$ gives \eqref{eq:xi-exact}.

For nonvanishing, the $j=3$ component has the unique lowest
antiholomorphic exponent: $n=-1$ contributes
$-\overline q^{1/40}$ to $\overline{g_3}$, while $T=-1$ contributes
$q^{1/8}$ to $\theta_3$.  No other $g_j$ has exponent $1/40$.
Hence the coefficient of $q^{1/8}\overline q^{1/40}$ in
\eqref{eq:dbar-exact} is nonzero.  Formula \eqref{eq:xi-exact}
contains the corresponding nonconstant factor
$\overline{\theta_3}$, so it is not holomorphic.
\end{proof}

\begin{proof}[Proof of Theorem~\ref{thm:main}]
The transformation statements are
Theorem~\ref{thm:zwegers-identification}; the differential statements
are Proposition~\ref{prop:shadow}.  For a weight-one harmonic Maass
form, the $\xi_1$ image is holomorphic of weight one.  Since
\eqref{eq:xi-exact} is not holomorphic, this completion is not
harmonic.  Thus this construction does not establish ordinary
weight-one mock modularity in the standard harmonic-Maass sense; see
\cite{jBF04} for the operator convention and
\cite{oBFOR17} for the harmonic-Maass framework.  The calculation
alone does not rule out the existence of
a different harmonic completion of $F$.
\end{proof}

%----------------------------------------------------------------------
\section{Comparison with the Maass cases}\label{sec:comparison}
%----------------------------------------------------------------------

The following table separates the proved differential calculation from the
structural comparison that motivates it.

\begin{center}
\small
\begin{tabular}{@{}lp{0.34\textwidth}p{0.34\textwidth}@{}}
\toprule
Feature & ADH/Cohen over $\QQ(\sqrt6)$ & $B$ over $\QQ(\sqrt5)$ \\
\midrule
Discriminant & $24$ & $20$ \\
Negative-vector norms & equal, automorph-related & $-5$ and $-20$ \\
Boundary corrections & cancel in the Maass construction & one vanishes, one survives \\
Differential image & Maass-type cancellation & mixed theta expression \\
Coefficient mechanism & Hecke-character descent & classical descent obstructed \\
Result & Maass waveform & nonharmonic indefinite completion \\
\bottomrule
\end{tabular}
\end{center}

In the ADH/Cohen case, the relevant equal-norm vectors are related by a
Pell automorphism, and the associated boundary terms combine to give
the Maass construction \cite{jADH88,jCo88,jZw12}.  Here the norm
inequality
\[
  P(c_1)=-5\neq-20=P(c_2)
\]
prevents any form-preserving automorphism from carrying one vector to
the other.  This explains why the classical symmetry is unavailable,
but it is not used as a converse theorem.  The analytic category is
decided instead by the exact differential image in
Section~\ref{sec:shadow}.

The transverse parity gives a second concrete contrast.  At $c_1$ the
primitive perpendicular direction is $(1,0)$, whose coordinate sum is
odd, and the shifted transverse theta is killed by
Lemma~\ref{lem:odd-theta}.  At $c_2$ the perpendicular direction is
$(3,-1)$, whose coordinate sum is even, and the four residue components
survive.  This concrete parity statement is sufficient for the present
example; we do not assert a general if-and-only-if criterion for every
shifted Pell lattice.

Zwegers' mock Maass theta functions \cite{jZw12} already show that
real-quadratic sign kernels can have nonholomorphic completions outside
the classical Maass-waveform examples.  The contribution here is the
exact placement of this cone series in that landscape: the
specific shift, the corrected boundary signs, the one-sided vanishing,
and the level-$40$ exponent normalization of the longitudinal theta
components.  In particular, the paper does not present
the general theory of mock Maass theta functions as new.

%----------------------------------------------------------------------
\section{Coefficient data, formalization scope, and open questions}
\label{sec:status}
%----------------------------------------------------------------------

\subsection{A coefficient witness}
The arithmetic function obtained from $B_N$ on integers
$M\equiv1\pmod{10}$ is not multiplicative.  The smallest displayed
witness from the companion paper is
\[
  B_1=1,\qquad B_3=-2,\qquad B_{34}=3,
\]
so that, for $11\cdot31=341$,
\[
  B_{(341-1)/10}=3\neq-2
  =B_{(11-1)/10}B_{(31-1)/10}.
\]
This fact is logically separate from Theorem~\ref{thm:main}: the
analytic classification uses the exact completion and its derivative,
not a contrapositive argument from nonmultiplicativity.

\subsection{Computational checks}
Direct enumeration of the cone and slab definitions agrees
coefficientwise well beyond the range needed for the examples in this
paper.  The accompanying script \texttt{slab\_check.py} performs an
independent exact-integer check for every $0\leq N\leq50$, prints the
two coefficients, and reports pass or fail for each $N$.  Its bounds
are derived from the cone inequalities rather than chosen as a search
radius.

\subsection{Coefficient statistics}
Through $N\leq10^5$: $\max_N|B_N|=7$, first attained at $N=27252$;
among the $19{,}617$ primes $p\equiv1\pmod{10}$ with
$(p-1)/10\leq10^5$, every value $B_{(p-1)/10}$ lies in
$\{-2,-1,+1,+2\}$, with magnitude $1$ for $13{,}097$ primes
($66.8\%$) and magnitude $2$ for $6520$ primes ($33.2\%$); there are
$4904$ indices at which
$10N+1$ is a norm but $B_N=0$, all composite.  Empirically
$|B_N|=O\bigl(d(10N+1)\bigr)$, a divisor-type bound.  The companion
paper explains the observed prime-magnitude split using an invariant
that combines an archimedean sector label with a congruence label.  Its
density statement is logically separate from the present analysis.

\subsection{Lean scope}
The Lean development checks the finite algebraic identities used for
the cone coefficients, the involution, the slab decomposition, the
Pell-coordinate formulas, the boundary-sign certificates, and a
finite residue-support surrogate, with no admitted lemmas in that
certificate layer.  The coefficient value $B_{34}=3$ is checked by
kernel evaluation and depends only on Lean's usual three axioms.
The analytic completion theorem from
Zwegers, the modular transformation law, and the analytic
differentiation in Proposition~\ref{prop:shadow} are not formalized.
The arithmetic foundations and their machine-checked scope are
described in the companion paper \cite{companion}.

\begin{remark}[Scope of the row-model provenance]
The corresponding Lean theorem proves the finite coefficient
convolution for \texttt{MKcoeff}.  Its declaration is
\texttt{mk\_factorization}, in namespace \texttt{Ch10} within
\texttt{QseriesFormalization}.  It is paired with
\texttt{coneDiffH\_two\_mul}, which identifies the even cone
coefficients with $B_N$, and \texttt{coneDiffH\_odd}, which proves
odd-level vanishing.  The Jacobi-triple-product evaluation of the
theta legs and the halved-variable Laurent-series restatement in
Theorem~\ref{thm:mk-factorization} are manuscript proofs, not
conclusions of that Lean endpoint.  No source-faithfulness theorem in
the repository currently identifies \texttt{MKcoeff} with the older
Chapter--10 object
\[
  \texttt{rhoEvenSourceOldRow\_missing}.
\]
Accordingly, no theorem about the full function $\Theta_{10}$ is
deduced here.
\end{remark}

\subsection{A family suggested by the example}
For squarefree $d\equiv5\pmod8$, the prime $2$ is inert in
$\QQ(\sqrt d)$.  Elementary mod-$8$ arithmetic on the Pell equations
$b^2-da'^2=1$ and $b^2-da'^2=4$ then forces alternating and constant
transverse characters, respectively, at the two boundary directions,
which suggests the one-sided pattern of Section~\ref{sec:shadow};
this parity computation has been checked numerically for
$d=5,13,29,37,53,61$.  The computation alone does not identify the
relevant boundary vectors for every $d$, so we state only a
conjecture.

\begin{conjecture}\label{conj:family}
Let $d\equiv5\pmod8$ be squarefree, and suppose a conductor-$2$
dissection over $\QQ(\sqrt d)$ produces a signature-$(1,1)$ shifted
lattice with boundary vectors $c_1(d)$, $c_2(d)$ of norms $-d$ and
$-4d$, positioned as in the $d=5$ model of
Section~\ref{sec:pell} (in particular
$\mathcal B(c_1(d),c_2(d))<0$ and the shifted lattice avoids the
boundary lines).  Then the correction at $c_1(d)$ vanishes by the odd
pairing of Lemma~\ref{lem:odd-theta}, the correction at $c_2(d)$
survives, and the residual indefinite theta function has a nonzero
mixed boundary derivative.
\end{conjecture}

\subsection{Open questions}
\begin{enumerate}
\item What is the minimal closed vector of characteristics containing
$\widehat F$ under $S$ and $T$, and what is its smallest congruence
subgroup realization?
\item Which other classical dissections produce one-sided or two-sided
boundary corrections?
\item Are the coefficients $B_N$ unbounded, and can one prove a
divisor-type upper bound from the norm interpretation?
\item For which squarefree $d$ can the conductor-$2$ boundary vectors
be determined uniformly enough to decide Conjecture~\ref{conj:family}?
\end{enumerate}

\subsection*{Acknowledgments}
The author thanks the readers who identified the separation needed
between the companion arithmetic foundations and the analytic content
of this paper.

\bibliographystyle{plain}
\bibliography{master}

\end{document}
